\documentclass[11pt]{article}
\usepackage[margin=0.8in]{geometry}
\usepackage{amsmath,microtype}
\newcommand{\guideheading}[1]{\par\medskip\noindent{\large\bfseries #1}\par\smallskip}
\begin{document}
\begin{center}
{\Large Guide: Using Merge Rules After an Orbit Shift}\par\medskip
Collatz proof project\quad---\quad September 5, 2026
\end{center}
\guideheading{What changed}
Earlier rules could replace certain hypothetical noncontracting seeds
with smaller ones. A noncontracting seed is one whose multiplicative
coefficient never falls below one. Moving forward along its orbit can
lose that property, so applying the old rule theorem directly to a
shifted value would leave a gap.

The new Lean proof keeps a numerical lower bound on every future
coefficient. A rule improves that bound. Several rules can restore it
to one, at which point the resulting seed is fully noncontracting again.
If that seed is smaller than the original, the original could not have
been the least positive noncontracting seed.

\guideheading{How the accounting works}
Represent the lower bound as $1/D$, with $D\geq1$. Moving forward by
$k$ steps from a fully noncontracting seed supplies the initial value
$D=C_k$, the coefficient accumulated during that shift. A merging rule
with coefficient boost $\rho$ updates the bound by
\[
 D\longmapsto\max(1,D/\rho).
\]
The rule must have an actual merging path and a noncontracting
replacement prefix. These hypotheses are checked; they are not assumed
for arbitrary paths. All calculations use exact integers or fractions.

\guideheading{A verified infinite family}
For every natural number $Q$, Lean proves the conditional implication
\[
447+549755813888Q\text{ is noncontracting}
\quad\Longrightarrow\quad
307+376572715308Q\text{ is noncontracting}.
\]
The second seed is always positive and smaller. At $Q=0$, the proof
shifts 447 forward 23 steps to 767, then applies three rules producing
511, 461 and 307. The intermediate deficits exceed one; the third rule
restores the bound to one. The theorem applies to every quotient, not
just these four displayed small seeds.

\guideheading{What the experiment says}
Of 2,336 finite seeds left unresolved by the earlier search, the new
search finds chain certificates for 1,012. For 182 of these, no single
rule suffices at any searched shift. Another 1,216 reach coefficient
contraction before a chain is found; 108 remain at the 64-step horizon.
These counts concern the actual finite seeds, not all numbers sharing
their initial parity patterns. No resource cap was reached. Independent
orbit replay checks every saved chain, but the full census is not a
Lean completeness theorem.

\guideheading{What still needs to be proved}
We need an argument that a suitable chain exists for every hypothetical
least positive noncontracting seed. Neither the finite search nor the
infinite family supplies that coverage. The nontrivial-cycle problem
also remains. Collatz is not proved.

\guideheading{Where to check the work}
\begingroup\raggedright
The formal proofs are in \texttt{lean/Collatz/MergeDeficit.lean}.
From \texttt{lean/}, run \texttt{lake build Collatz.MergeDeficit}.
From the root, run \texttt{python3 analysis/merge\_chain\_search.py} to
reproduce the census, or \texttt{python3 -m unittest discover -s analysis
-p test\_merge\_chain\_search.py} for independent checks.
The technical paper gives the exact hypotheses and family arithmetic.
\par\endgroup
\end{document}
